% TeX program = pdflatex
\documentclass[11pt,a4paper]{amsart}

% Shared preamble for course notes and exercise sets.

\usepackage[T1]{fontenc}
\usepackage{lmodern}
\usepackage[english]{babel}

\usepackage[a4paper,margin=30mm]{geometry}
\usepackage{microtype}
\usepackage{graphicx}
\usepackage{enumitem}

\usepackage{amsmath}
\usepackage{amssymb}
\usepackage{amsthm}
\usepackage{mathtools}
\usepackage{aliascnt}

\usepackage[hidelinks,pdfusetitle]{hyperref}

\numberwithin{equation}{section}
\setlist{itemsep=0.25em,topsep=0.5em}

\theoremstyle{plain}
\newtheorem{theorem}{Theorem}[section]
\newaliascnt{lemma}{theorem}
\newtheorem{lemma}[lemma]{Lemma}
\aliascntresetthe{lemma}
\newaliascnt{proposition}{theorem}
\newtheorem{proposition}[proposition]{Proposition}
\aliascntresetthe{proposition}
\newaliascnt{corollary}{theorem}
\newtheorem{corollary}[corollary]{Corollary}
\aliascntresetthe{corollary}

\theoremstyle{definition}
\newaliascnt{definition}{theorem}
\newtheorem{definition}[definition]{Definition}
\aliascntresetthe{definition}
\newaliascnt{example}{theorem}
\newtheorem{example}[example]{Example}
\aliascntresetthe{example}
\newaliascnt{problem}{theorem}
\newtheorem{problem}[problem]{Problem}
\aliascntresetthe{problem}

\theoremstyle{remark}
\newaliascnt{remark}{theorem}
\newtheorem{remark}[remark]{Remark}
\aliascntresetthe{remark}

\usepackage[nameinlink,noabbrev]{cleveref}

\crefname{theorem}{theorem}{theorems}
\crefname{lemma}{lemma}{lemmas}
\crefname{proposition}{proposition}{propositions}
\crefname{corollary}{corollary}{corollaries}
\crefname{definition}{definition}{definitions}
\crefname{example}{example}{examples}
\crefname{problem}{problem}{problems}
\crefname{remark}{remark}{remarks}

\DeclareMathOperator{\im}{im}
\DeclareMathOperator{\rank}{rank}
\DeclarePairedDelimiter{\abs}{\lvert}{\rvert}
\DeclarePairedDelimiter{\norm}{\lVert}{\rVert}
\DeclarePairedDelimiter{\set}{\lbrace}{\rbrace}


\title{Introduction to Logic I}
\author{Joona Piirainen}
\date{Autumn 2026}

\begin{document}

\maketitle

\begin{abstract}
  Notes for the University of Helsinki course MAT21014, taught in Finnish
  from 03.09.2026--17.10.2026.
\end{abstract}

\tableofcontents

\section{Course information}

The official course information is available on the
\href{https://studies.helsinki.fi/kurssit/opintojakso/otm-836bfa0a-198b-4cab-b040-6c0bd165b235/MAT21014}{University
of Helsinki course page}.

\section{Lecture notes}

\subsection{03.09.2026}

\subsubsection{Propositional formulas}

\begin{definition}
  The set of \emph{propositional formulas} is defined inductively as follows:
  \begin{enumerate}
    \item Every propositional symbol \(p_n\), where \(n \in \NN\), is a
      propositional formula.
    \item If \(A\) and \(B\) are propositional formulas, then
      \[
        \neg A, \qquad
        (A \land B), \qquad
        (A \lor B), \qquad
        (A \rightarrow B), \qquad
        (A \leftrightarrow B)
      \]
      are propositional formulas.
    \item No other strings are propositional formulas.
  \end{enumerate}
\end{definition}

\subsubsection{Subformulas}

\begin{definition}[Immediate subformulas]
  The immediate subformulas of a propositional formula are defined by its
  outermost connective:
  \begin{enumerate}
    \item A propositional symbol \(p_n\) has no immediate subformulas.
    \item The only immediate subformula of \(\neg A\) is \(A\).
    \item The immediate subformulas of \((A \mathbin{\circ} B)\) are \(A\)
      and \(B\), where
      \[
        \circ \in \{\land, \lor, \rightarrow, \leftrightarrow\}.
      \]
  \end{enumerate}
\end{definition}

\begin{definition}[Subformulas]
  A formula \(B\) is a \emph{subformula} of a formula \(A\) if \(B=A\), or
  if there are formulas
  \[
    A=A_0,A_1,\ldots,A_k=B
  \]
  such that \(k \geq 1\) and \(A_{i+1}\) is an immediate subformula of
  \(A_i\) for every \(0 \leq i<k\).
\end{definition}

\subsubsection{Structural induction}

\begin{theorem}[Structural induction]
  Let \(P(A)\) be a statement about propositional formulas. Suppose that
  \begin{enumerate}
    \item \(P(p_n)\) holds for every \(n \in \NN\);
    \item if \(P(A)\) holds, then \(P(\neg A)\) holds; and
    \item if \(P(A)\) and \(P(B)\) hold, then
      \(P((A \mathbin{\circ} B))\) holds for every
      \[
        \circ \in \{\land, \lor, \rightarrow, \leftrightarrow\}.
      \]
  \end{enumerate}
  Then \(P(A)\) holds for every propositional formula \(A\).
\end{theorem}

\begin{example}[Structural induction]
  For a formula \(A\), let \(v(A)\) be the number of occurrences of
  propositional symbols in \(A\), and let \(b(A)\) be the number of
  occurrences of binary connectives in \(A\). We prove that
  \[
    v(A)=b(A)+1
  \]
  for every propositional formula \(A\).

  If \(A=p_n\), then \(v(A)=1\) and \(b(A)=0\), so the claim holds. Suppose
  next that the claim holds for a formula \(A\). Negation introduces neither
  a propositional symbol nor a binary connective, and hence
  \[
    v(\neg A)=v(A)=b(A)+1=b(\neg A)+1.
  \]
  Finally, suppose that the claim holds for formulas \(A\) and \(B\), and let
  \(\circ\) be any binary connective. Then
  \begin{align*}
    v((A \mathbin{\circ} B))
    &=v(A)+v(B) \\
    &=(b(A)+1)+(b(B)+1) \\
    &=b((A \mathbin{\circ} B))+1.
  \end{align*}
  The result follows by structural induction. \qed
\end{example}

\subsection{07.09.2026}

\subsubsection{Structural induction}

\begin{example}[Matching parentheses]
  For a propositional formula \(A\), let \(l(A)\) and \(r(A)\) denote the
  numbers of occurrences of left and right parentheses in \(A\), respectively.
  We prove by structural induction that
  \[
    l(A)=r(A)
  \]
  for every propositional formula \(A\).

  If \(A=p_n\), then \(l(A)=0=r(A)\). Suppose next that the claim holds for
  a formula \(A\). Negation introduces no parentheses, so
  \[
    l(\neg A)=l(A)=r(A)=r(\neg A).
  \]
  Finally, suppose that the claim holds for formulas \(A\) and \(B\), and let
  \(\circ\) be any binary connective. The formula \((A \mathbin{\circ} B)\)
  introduces one left and one right parenthesis. Hence
  \begin{align*}
    l((A \mathbin{\circ} B))
    &=l(A)+l(B)+1 \\
    &=r(A)+r(B)+1 \\
    &=r((A \mathbin{\circ} B)).
  \end{align*}
  The result follows by structural induction. \qed
\end{example}

\subsubsection{Closures}

\begin{definition}[Closure]
  Let \(X\) be a set, let \(B \subseteq X\), and, for \(1 \leq i \leq k\), let
  \[
    f_i\colon X^{n_i} \longrightarrow X
  \]
  be a function. A subset \(Y \subseteq X\) is \emph{closed under \(f_i\)} if
  \[
    f_i(y_1,\ldots,y_{n_i}) \in Y
  \]
  whenever \(y_1,\ldots,y_{n_i} \in Y\). The \emph{closure of \(B\) under
  \(f_1,\ldots,f_k\)}, denoted by
  \[
    \operatorname{cl}(B;f_1,\ldots,f_k),
  \]
  is the smallest subset of \(X\) that contains \(B\) and is closed under
  every \(f_i\).
\end{definition}

Equivalently, the closure can be written as the intersection
\[
  \operatorname{cl}(B;f_1,\ldots,f_k)
  =\bigcap\{Y \subseteq X : B \subseteq Y
  \text{ and \(Y\) is closed under every \(f_i\)}\}.
\]

\begin{example}
  Let \(X=\ZZ\), \(B=\{1\}\), and define functions \(f,g\colon\ZZ\to\ZZ\) by
  \[
    f(x)=x+1
    \qquad\text{and}\qquad
    g(x)=-x.
  \]
  Then
  \[
    \operatorname{cl}(\{1\};f,g)=\ZZ.
  \]
  Indeed, every subset of \(\ZZ\) that contains \(1\) and is closed under
  \(g\) must also contain \(-1\). Closure under \(f\) then gives \(0\), and
  repeated applications of \(f\) give every positive integer. Applying \(g\)
  to these integers gives every negative integer. Thus any such closed set is
  all of \(\ZZ\), while \(\ZZ\) itself clearly contains \(1\) and is closed
  under both functions.
\end{example}

\begin{corollary}[Induction on a closure]
  Set
  \[
    C=\operatorname{cl}(B;f_1,\ldots,f_k).
  \]
  To prove that every element of \(C\) satisfies a property \(P\), it suffices
  to prove that
  \begin{enumerate}
    \item every element of \(B\) satisfies \(P\); and
    \item for each \(i\), if \(x_1,\ldots,x_{n_i}\) satisfy \(P\), then
      \(f_i(x_1,\ldots,x_{n_i})\) also satisfies \(P\).
  \end{enumerate}
\end{corollary}

\begin{proof}
  Let \(Y\) be the set of all elements of
  \(\operatorname{cl}(B;f_1,\ldots,f_k)\) that satisfy \(P\). The two
  assumptions imply that \(B \subseteq Y\) and that \(Y\) is closed under
  every \(f_i\). Since the closure is the smallest such set,
  \[
    \operatorname{cl}(B;f_1,\ldots,f_k) \subseteq Y.
  \]
  Hence every element of the closure satisfies \(P\).
\end{proof}

Let \(\Sigma\) be the alphabet of propositional logic, let \(\Sigma^*\) be the
set of all finite strings over \(\Sigma\), and let
\[
  V=\{p_n:n\in\NN\}
\]
be the set of propositional symbols. The logical connectives determine the
following functions on strings:
\[
  f_{\neg}(A)=\neg A
\]
and
\[
  f_{\circ}(A,B)=(A \mathbin{\circ} B),
  \qquad
  \circ\in\{\land,\lor,\rightarrow,\leftrightarrow\}.
\]
The set of propositional formulas is precisely
\[
  \operatorname{cl}
  (V;f_{\neg},f_{\land},f_{\lor},f_{\rightarrow},f_{\leftrightarrow}).
\]
In particular, if \(S \subseteq \Sigma^*\) contains every propositional symbol
and is closed under all functions associated with the logical connectives,
then \(S\) contains every propositional formula.

\subsubsection{Semantics of propositional logic}

\begin{definition}[Truth assignment]
  A \emph{truth assignment} is a function
  \[
    v\colon V\longrightarrow\{0,1\},
  \]
  where \(V\) is the set of propositional symbols. The values \(0\) and \(1\)
  represent \textbf{false} and \textbf{true}, respectively.
\end{definition}

The truth values of the logical connectives are defined as follows.

\begin{definition}[Negation]
  The negation \(\neg A\) is true exactly when \(A\) is false.
  \[
    \begin{array}{c|c}
      A & \neg A \\
      \hline
      0 & 1 \\
      1 & 0
    \end{array}
  \]
\end{definition}

\begin{definition}[Conjunction]
  The conjunction \((A\land B)\) is true exactly when both \(A\) and \(B\) are
  true.
  \[
    \begin{array}{cc|c}
      A & B & (A\land B) \\
      \hline
      0 & 0 & 0 \\
      0 & 1 & 0 \\
      1 & 0 & 0 \\
      1 & 1 & 1
    \end{array}
  \]
\end{definition}

\begin{definition}[Disjunction]
  The disjunction \((A\lor B)\) is true exactly when at least one of \(A\) and
  \(B\) is true.
  \[
    \begin{array}{cc|c}
      A & B & (A\lor B) \\
      \hline
      0 & 0 & 0 \\
      0 & 1 & 1 \\
      1 & 0 & 1 \\
      1 & 1 & 1
    \end{array}
  \]
\end{definition}

\begin{definition}[Implication]
  The implication \((A\rightarrow B)\) is false exactly when \(A\) is true and
  \(B\) is false. It is important to note that implication is \textbf{not}
  symmetric.
  \[
    \begin{array}{cc|c}
      A & B & (A\rightarrow B) \\
      \hline
      0 & 0 & 1 \\
      0 & 1 & 1 \\
      1 & 0 & 0 \\
      1 & 1 & 1
    \end{array}
  \]
\end{definition}

\begin{definition}[Equivalence]
  The equivalence \((A\leftrightarrow B)\) is true exactly when \(A\) and \(B\)
  have the same truth value.
  \[
    \begin{array}{cc|c}
      A & B & (A\leftrightarrow B) \\
      \hline
      0 & 0 & 1 \\
      0 & 1 & 0 \\
      1 & 0 & 0 \\
      1 & 1 & 1
    \end{array}
  \]
\end{definition}

\begin{definition}[Truth value]
  Let \(A\) be a propositional formula and let \(v\) be a truth assignment.
  The \emph{truth value} \(v(A)\) of \(A\) is defined by the following rules:
  \begin{enumerate}
    \item If \(A\) is a propositional symbol \(p_n\), then \(v(A)\) is already
      defined.
    \item If \(A=\neg B\) and \(v(B)\) is defined, then
      \[
        v(A)=
        \begin{cases}
          0, & \text{if }v(B)=1, \\
          1, & \text{if }v(B)=0.
        \end{cases}
      \]
      Another way to write this is \(v(A) = 1 - v(B)\).
    \item If \(A=(B\land C)\) and \(v(B)\) and \(v(C)\) are defined, then
      \(v(A)=1\) exactly when \(v(B)=v(C)=1\), and \(v(A)=0\) otherwise.
      Another way to write this is \(v(A)v(B)\).
    \item If \(A=(B\lor C)\) and \(v(B)\) and \(v(C)\) are defined, then
      \(v(A)=1\) when at least one of \(v(B)=1\) and \(v(C)=1\) holds, and
      \(v(A)=0\) when \(v(B)=v(C)=0\).
    \item If \(A=(B\rightarrow C)\) and \(v(B)\) and \(v(C)\) are defined,
      then \(v(A)=0\) when \(v(B)=1\) and \(v(C)=0\), and \(v(A)=1\) in all
      other cases.
    \item If \(A=(B\leftrightarrow C)\) and \(v(B)\) and \(v(C)\) are
      defined, then \(v(A)=1\) iff \(v(B)=v(C)\).
  \end{enumerate}
\end{definition}

\begin{definition}[Classification of propositional formulas]
  A propositional formula \(A\) is
  \begin{enumerate}
    \item a \emph{tautology} if \(v(A)=1\) for every truth assignment \(v\);
    \item a \emph{contradiction} if \(v(A)=0\) for every truth assignment
      \(v\);
    \item \emph{contingent} if there are truth assignments \(v\) and \(w\)
      such that \(v(A)=1\) and \(w(A)=0\);
    \item \emph{satisfiable} if \(v(A)=1\) for at least one truth assignment
      \(v\); and
    \item \emph{falsifiable} if \(v(A)=0\) for at least one truth assignment
      \(v\).
  \end{enumerate}
\end{definition}

\begin{example}[De Morgan's laws]
  Consider the formulas
  \[
    (\neg(A\land B)\leftrightarrow(\neg A\lor\neg B))
    \qquad\text{and}\qquad
    (\neg(A\lor B)\leftrightarrow(\neg A\land\neg B)).
  \]
  Their truth values can be compared using the following table:
  \[
    \begin{array}{cc|cc|cc}
      A & B & \neg(A\land B) & (\neg A\lor\neg B)
      & \neg(A\lor B) & (\neg A\land\neg B) \\
      \hline
      0 & 0 & 1 & 1 & 1 & 1 \\
      0 & 1 & 1 & 1 & 0 & 0 \\
      1 & 0 & 1 & 1 & 0 & 0 \\
      1 & 1 & 0 & 0 & 0 & 0
    \end{array}
  \]
  The two sides of each equivalence have the same truth value under every
  truth assignment. Consequently, both equivalences are tautologies. These
  equivalences are known as \emph{De Morgan's laws}.
\end{example}

\subsection{10.09.2026}

\subsubsection{Logical implication and equivalence}

\begin{definition}[Logical implication]
  A formula \(A\) \emph{logically implies} a formula \(B\), written
  \[
    A \Rightarrow B,
  \]
  if, for every truth assignment \(v\), \(v(A)=1\) implies \(v(B)=1\).
  Equivalently, \(A \Rightarrow B\) exactly when \((A\rightarrow B)\) is a
  tautology.
\end{definition}

\begin{definition}[Logical equivalence]
  Formulas \(A\) and \(B\) are \emph{logically equivalent}, written
  \[
    A \Leftrightarrow B,
  \]
  if \(v(A)=v(B)\) for every truth assignment \(v\). Equivalently,
  \(A \Leftrightarrow B\) exactly when \((A\leftrightarrow B)\) is a
  tautology.
\end{definition}

The symbols \(\Rightarrow\) and \(\Leftrightarrow\) are used in the
metalanguage to state relations between formulas; unlike \(\rightarrow\) and
\(\leftrightarrow\), they are not connectives used to construct propositional
formulas.

\begin{definition}[Literal]
  A \emph{literal} is either a propositional symbol \(p_n\) or its negation
  \(\neg p_n\). The former is a \emph{positive literal}, and the latter is a
  \emph{negative literal}.
\end{definition}

\begin{definition}[Disjunctive normal form]
  A formula is in \emph{disjunctive normal form} (DNF) if it is a disjunction
  of conjunctions of literals. In other words, it has the form
  \[
    \bigvee_{i=1}^m\left(\bigwedge_{j=1}^{k_i}L_{ij}\right),
  \]
  where \(m\geq 1\), each \(k_i\geq 1\), and every \(L_{ij}\) is a literal.
\end{definition}

\begin{definition}[Conjunctive normal form]
  A formula is in \emph{conjunctive normal form} (CNF) if it is a conjunction
  of disjunctions of literals. In other words, it has the form
  \[
    \bigwedge_{i=1}^m\left(\bigvee_{j=1}^{k_i}L_{ij}\right),
  \]
  where \(m\geq 1\), each \(k_i\geq 1\), and every \(L_{ij}\) is a literal.
\end{definition}

For a literal \(L\), let \(\overline L\) denote its \emph{complementary
literal}, defined by
\[
  \overline{p_n}=\neg p_n
  \qquad\text{and}\qquad
  \overline{\neg p_n}=p_n.
\]
De Morgan's laws show that DNF and CNF are dual: negating a DNF formula gives
an equivalent CNF formula, and negating a CNF formula gives an equivalent DNF
formula. More precisely,
\[
  \neg\bigvee_{i=1}^m\left(\bigwedge_{j=1}^{k_i}L_{ij}\right)
  \Leftrightarrow
  \bigwedge_{i=1}^m\left(\bigvee_{j=1}^{k_i}\overline{L_{ij}}\right),
\]
and the analogous equivalence holds with \(\land\) and \(\lor\) interchanged.

\begin{definition}[Boolean function]
  An \(n\)-ary \emph{truth function} is a
  function
  \[
    f\colon\{0,1\}^n\longrightarrow\{0,1\}.
  \]
  A propositional formula \(A\) whose propositional symbols are among
  \(p_1,\ldots,p_n\) determines the truth function
  \[
    f_A(x_1,\ldots,x_n)=v(A),
  \]
  where \(v(p_i)=x_i\) for every \(1\leq i\leq n\).
\end{definition}

\begin{theorem}[Normal form theorem]
  Let \(f\) be an \(n\)-ary truth function, where \(n\geq 1\). Then there are
  formulas \(A_D\) in DNF and \(A_C\) in CNF, containing only the
  propositional symbols \(p_1,\ldots,p_n\), such that
  \[
    f_{A_D}=f_{A_C}=f.
  \]
  Consequently, every propositional formula is logically equivalent to a
  formula in DNF and to a formula in CNF.
\end{theorem}

\begin{proof}
  We construct the formulas directly from the truth table of \(f\).

  To construct \(A_D\), consider each row
  \(a=(a_1,\ldots,a_n)\in\{0,1\}^n\) on which \(f(a)=1\), and define the
  formula
  \[
    M_a=L_1^a\land\cdots\land L_n^a,
    \qquad
    L_i^a=
    \begin{cases}
      p_i, & \text{if }a_i=1,\\
      \neg p_i, & \text{if }a_i=0.
    \end{cases}
  \]
  Thus \(M_a\) is true under a truth assignment \(v\) exactly when
  \((v(p_1),\ldots,v(p_n))=a\).

  If there is at least one such row, let
  \[
    A_D=\bigvee_{\substack{a\in\{0,1\}^n\\f(a)=1}}M_a.
  \]
  Under an assignment corresponding to \(b\in\{0,1\}^n\), the formula \(A_D\)
  is true exactly when the disjunction contains \(M_b\), which happens exactly
  when \(f(b)=1\). If \(f\) is constantly false, take
  \(A_D=(p_1\land\neg p_1)\). Thus \(f_{A_D}=f\).

  Dually, to construct \(A_C\), consider each row \(a\) on which \(f(a)=0\)
  and define
  \[
    C_a=K_1^a\lor\cdots\lor K_n^a,
    \qquad
    K_i^a=
    \begin{cases}
      p_i, & \text{if }a_i=0,\\
      \neg p_i, & \text{if }a_i=1.
    \end{cases}
  \]
  The clause \(C_a\) is false exactly on the row \(a\). If there is at least
  one such row, let
  \[
    A_C=\bigwedge_{\substack{a\in\{0,1\}^n\\f(a)=0}}C_a.
  \]
  This conjunction is true on a row \(b\) exactly when \(f(b)=1\). If \(f\)
  is constantly true, take \(A_C=(p_1\lor\neg p_1)\). Thus
  \(f_{A_C}=f\).

  All finite conjunctions and disjunctions above may be given any fixed binary
  parenthesization. Finally, for any propositional formula, list its distinct
  propositional symbols as \(q_1,\ldots,q_n\) and apply the same construction
  with \(q_i\) in place of \(p_i\). This gives logically equivalent formulas
  in DNF and CNF.
\end{proof}

\begin{example}[Constructing normal forms from a truth function]
  Consider the binary truth function whose truth table is
  \[
    \begin{array}{cc|c}
      x_1 & x_2 & f(x_1,x_2)\\
      \hline
      0 & 0 & 0\\
      0 & 1 & 1\\
      1 & 0 & 1\\
      1 & 1 & 0
    \end{array}
  \]
  The rows where \(f=1\) are \((0,1)\) and \((1,0)\). Replace a \(1\) in a
  row by \(p_i\), replace a \(0\) by \(\neg p_i\), and join the literals in
  each row with \(\land\). This gives the formulas
  \[
    M_{(0,1)}=(\neg p_1\land p_2)
    \qquad\text{and}\qquad
    M_{(1,0)}=(p_1\land\neg p_2).
  \]
  Finally, join the formulas for the selected rows with \(\lor\):
  \[
    A=((\neg p_1\land p_2)\lor(p_1\land\neg p_2)).
  \]
  This is a DNF formula with \(f_A=f\).

  To construct a CNF formula, use instead the rows \((0,0)\) and \((1,1)\),
  where \(f=0\). The clauses that are false on exactly these rows are
  \[
    C_{(0,0)}=(p_1\lor p_2)
    \qquad\text{and}\qquad
    C_{(1,1)}=(\neg p_1\lor\neg p_2).
  \]
  Their conjunction
  \[
    B=((p_1\lor p_2)\land(\neg p_1\lor\neg p_2))
  \]
  is in CNF and satisfies \(f_B=f\). Thus \(A\Leftrightarrow B\).
\end{example}

\subsection{14.09.2026}

\subsubsection{Universal sets of connectives}

\begin{definition}[Universal set of connectives]
  A set \(C\) of logical connectives is \emph{universal}, or
  \emph{functionally complete}, if, for every \(n\geq 1\) and every truth
  function
  \[
    f\colon\{0,1\}^n\longrightarrow\{0,1\},
  \]
  there is a formula \(A\), containing only the propositional symbols
  \(p_1,\ldots,p_n\) and connectives from \(C\), such that \(f_A=f\).
\end{definition}

\begin{theorem}
  The set \(\{\neg,\land\}\) is universal.
\end{theorem}

\begin{proof}
  Let \(f\colon\{0,1\}^n\to\{0,1\}\) be a truth function. The DNF
  construction gives a DNF formula \(A\) such that \(f_A=f\). We prove by
  structural induction that every formula built from literals using
  \(\land\) and \(\lor\) is logically equivalent to a formula containing only
  \(\neg\) and \(\land\).

  A literal is either \(p_i\) or \(\neg p_i\), so the claim holds in the base
  case. Suppose that the induction hypothesis gives formulas \(B^*\) and
  \(C^*\), containing only \(\neg\) and \(\land\), such that
  \(B^*\Leftrightarrow B\) and \(C^*\Leftrightarrow C\). Then
  \[
    (B\land C)\Leftrightarrow(B^*\land C^*)
  \]
  and, by De Morgan's law,
  \[
    (B\lor C)\Leftrightarrow\neg(\neg B^*\land\neg C^*).
  \]
  Both formulas on the right use only \(\neg\) and \(\land\). Structural
  induction therefore gives a formula \(A^*\), containing only these
  connectives, such that \(A^*\Leftrightarrow A\). Hence
  \(f_{A^*}=f_A=f\), so \(\{\neg,\land\}\) is universal.
\end{proof}

\begin{theorem}
  The set \(\{\land,\lor\}\) is not universal.
\end{theorem}

\begin{proof}
  We prove by structural induction that every formula \(A\) containing only
  \(\land\) and \(\lor\) is false under the truth assignment \(v\) for which
  \(v(p_i)=0\) for every propositional symbol \(p_i\).

  If \(A=p_i\), then \(v(A)=0\). Suppose that the claim holds for formulas
  \(B\) and \(C\). Then
  \[
    v(B\land C)=0
    \qquad\text{and}\qquad
    v(B\lor C)=0,
  \]
  so the claim also holds for formulas formed using either allowed connective.
  It therefore holds for every formula containing only \(\land\) and \(\lor\).

  The unary truth function \(f(x)=1-x\) has value \(f(0)=1\), so no such
  formula can represent it. Hence \(\{\land,\lor\}\) is not universal.
\end{proof}

\subsection{21.09.2026}

\subsubsection{Clauses and clause sets}

\begin{definition}[Clause]
  A \emph{clause} is a finite set of literals. It represents the disjunction
  of its members. For a truth assignment \(v\), define
  \[
    v(C)=1 \quad\text{if and only if}\quad
    v(q)=1\text{ for some }q\in C,
  \]
  and set \(v(C)=0\) otherwise. In particular, the \emph{empty clause}
  \(\emptyset\) is false under every truth assignment.
\end{definition}

\begin{definition}[Clause set]
  A \emph{clause set} \(\mathcal C\) is a set of clauses. Define
  \[
    v(\mathcal C)=1 \quad\text{if and only if}\quad
    v(C)=1\text{ for every }C\in\mathcal C,
  \]
  and set \(v(\mathcal C)=0\) otherwise. Thus a finite clause set represents
  the conjunction of its clauses. It is \emph{satisfiable} if some truth
  assignment makes it true, and \emph{unsatisfiable} otherwise.
\end{definition}

\begin{remark}
  The empty clause and the empty clause set have different meanings:
  the empty clause is always false, whereas the empty clause set is always
  true (the condition on all its members holds vacuously). In particular,
  \(\{\emptyset\}\) is an unsatisfiable clause set.
  A clause containing both \(p_i\) and \(\neg p_i\) is always true.
\end{remark}

\begin{example}[Encoding CNF as a clause set]
  Since \((p_0\rightarrow p_1)\Leftrightarrow(\neg p_0\lor p_1)\), the formula
  \[
    p_0\land(p_0\rightarrow p_1)
  \]
  is represented by \(\mathcal C=\{\{p_0\},\{\neg p_0,p_1\}\}\).
  The literals within a clause are joined by disjunction, and the clauses
  within a clause set are joined by conjunction. Order and repetitions do
  not matter in these set representations.
\end{example}

\subsubsection{The resolution rule and resolution proofs}

\begin{definition}[Resolution rule]
  The \emph{resolution rule} is
  \[
    \frac{C_1\cup\{p_i\}\qquad C_2\cup\{\neg p_i\}}{C_1\cup C_2}.
  \]
  The conclusion is called a \emph{resolvent}, and \(p_i\) is the
  \emph{pivot}. In the usual application to clauses \(D,E\) with
  \(p_i\in D\) and \(\neg p_i\in E\), take
  \(C_1=D\setminus\{p_i\}\) and \(C_2=E\setminus\{\neg p_i\}\).
  Only the selected complementary pair is removed.
\end{definition}

\begin{example}[Resolution steps]
  Modus ponens is the resolution step
  \[
    \frac{\{p_0\}\qquad\{\neg p_0,p_1\}}{\{p_1\}}.
  \]
  From the same two clauses one may obtain different resolvents by choosing
  different pivots. For example,
  \[
    \frac{\{p_0,p_1\}\qquad\{\neg p_0,\neg p_1\}}
    {\{p_0,\neg p_0\}}
    \qquad\text{or}\qquad
    \frac{\{p_0,p_1\}\qquad\{\neg p_0,\neg p_1\}}
    {\{p_1,\neg p_1\}}.
  \]
  Both conclusions are tautological clauses. Removing both complementary
  pairs at once to obtain \(\emptyset\) would be invalid: the premises are
  simultaneously true when \(v(p_0)=1\) and \(v(p_1)=0\).
\end{example}

\begin{definition}[Resolution proof]
  A \emph{resolution proof} of a clause \(C\) from a clause set
  \(\mathcal C\) is a finite sequence \(C_0,\ldots,C_m\), where \(m\geq0\),
  such that \(C_m=C\) and, for every \(i\leq m\), either
  \begin{enumerate}
    \item \(C_i\in\mathcal C\), or
    \item \(C_i\) follows by resolution from \(C_k,C_l\) for some \(k,l<i\).
  \end{enumerate}
  We write \(\mathcal C\vdash_R C\) if such a proof exists.
  A proof of \(\emptyset\) is called a \emph{resolution refutation}.
\end{definition}

\subsubsection{Soundness of resolution}

\begin{lemma}[Soundness of the resolution rule]
  If \(v(C_1\cup\{p_i\})=v(C_2\cup\{\neg p_i\})=1\), then
  \(v(C_1\cup C_2)=1\).
\end{lemma}

\begin{proof}
  If \(v(p_i)=1\), then \(v(\neg p_i)=0\). Since the second premise is
  true, some literal of \(C_2\) must be true. Hence \(v(C_1\cup C_2)=1\).
  If \(v(p_i)=0\), truth of the first premise similarly implies that some
  literal of \(C_1\) is true, so again \(v(C_1\cup C_2)=1\).
\end{proof}

\begin{theorem}[Soundness]
  If \(\mathcal C\vdash_R C\), then every truth assignment satisfying
  \(\mathcal C\) satisfies \(C\). Thus every resolution conclusion is a
  logical consequence of its assumptions.
\end{theorem}

\begin{proof}
  Let \(C_0,\ldots,C_m=C\) be a resolution proof from \(\mathcal C\), and
  let \(v(\mathcal C)=1\). We prove \(v(C_i)=1\) by strong induction on \(i\).
  For \(i=0\), there are no earlier lines, so \(C_0\in\mathcal C\) and
  \(v(C_0)=1\). Suppose \(i>0\) and \(v(C_j)=1\) for every \(j<i\).
  If \(C_i\in\mathcal C\), it is true by assumption. Otherwise it is a
  resolvent of \(C_k,C_l\) with \(k,l<i\). Both premises are true by the
  induction hypothesis, and the soundness lemma gives \(v(C_i)=1\).
  In particular, \(v(C_m)=v(C)=1\).
\end{proof}

\begin{remark}[Completeness versus refutation completeness]
  Resolution does not derive every logical consequence: from
  \(\{\{p_0\}\}\) one cannot derive \(\{p_0,p_1\}\), although
  \(p_0\Rightarrow(p_0\lor p_1)\). Indeed, the only available clause is
  \(\{p_0\}\), and no resolution step is possible without a complementary
  literal. Resolution is nevertheless \emph{refutation complete}: an
  unsatisfiable clause set admits a derivation of the empty clause.
\end{remark}

\subsubsection{Proving an implication by refutation}

To prove \(A\Rightarrow B\), it suffices to show that \(A\land\neg B\)
is unsatisfiable: a truth assignment falsifying the implication would make
both \(A\) and \(\neg B\) true. Convert \(A\land\neg B\) into CNF,
encode it as a clause set, and derive \(\emptyset\).

\begin{example}
  We prove
  \[
    (p_0\land(p_0\rightarrow p_1))\Rightarrow(p_1\lor p_2).
  \]
  Negating the conclusion gives
  \(\neg(p_1\lor p_2)\Leftrightarrow(\neg p_1\land\neg p_2)\).
  Hence the premises together with the negated conclusion give
  \[
    \mathcal C=
    \bigl\{\{p_0\},\{\neg p_0,p_1\},\{\neg p_1\},\{\neg p_2\}\bigr\}.
  \]
  A resolution refutation is
  \[
    \begin{array}{rll}
      0.&\{p_0\}&\text{assumption},\\
      1.&\{\neg p_0,p_1\}&\text{assumption},\\
      2.&\{\neg p_1\}&\text{assumption},\\
      3.&\{\neg p_2\}&\text{assumption},\\
      4.&\{p_1\}&\text{resolution of lines 0 and 1},\\
      5.&\emptyset&\text{resolution of lines 4 and 2}.
    \end{array}
  \]
  Thus \(\mathcal C\vdash_R\emptyset\), so soundness implies that
  \(\mathcal C\) is unsatisfiable and proves the implication. Line 3 need
  not be used: the premises already imply \(p_1\).
\end{example}

\subsubsection{Refutation completeness}

\begin{theorem}[Finite refutation completeness]
  If \(\mathcal C\) is a finite unsatisfiable clause set, then
  \(\mathcal C\vdash_R\emptyset\).
\end{theorem}

\begin{proof}
  We use strong induction on the number \(n\) of distinct propositional
  symbols occurring in \(\mathcal C\). This number is finite because both
  \(\mathcal C\) and each of its clauses are finite.

  If \(n=0\), the only possible clause is \(\emptyset\).
  Since the empty clause set is satisfiable, unsatisfiability forces
  \(\emptyset\in\mathcal C\). The one-line proof \((\emptyset)\) suffices.

  Suppose \(n>0\) and the result holds for fewer than \(n\) symbols.
  Choose a symbol \(p_i\) occurring in \(\mathcal C\), and partition the
  clause set as in the book:
  \[
    \begin{aligned}
      \mathcal C_{00}&=\{C\in\mathcal C:p_i\notin C,\ \neg p_i\notin C\},\\
      \mathcal C_{01}&=\{C\in\mathcal C:p_i\notin C,\ \neg p_i\in C\},\\
      \mathcal C_{10}&=\{C\in\mathcal C:p_i\in C,\ \neg p_i\notin C\},\\
      \mathcal C_{11}&=\{C\in\mathcal C:p_i\in C,\ \neg p_i\in C\}.
    \end{aligned}
  \]
  Every clause in \(\mathcal C_{11}\) is tautological. Form all resolvents
  between the negative and positive parts on the pivot \(p_i\):
  \[
    \mathcal C_{01}\times_R\mathcal C_{10}
    =\bigl\{(C_1\setminus\{\neg p_i\})\cup(C_2\setminus\{p_i\}):
    C_1\in\mathcal C_{01},\ C_2\in\mathcal C_{10}\bigr\}.
  \]
  Define
  \[
    \mathcal C'=\mathcal C_{00}\cup
    (\mathcal C_{01}\times_R\mathcal C_{10}).
  \]
  No clause of \(\mathcal C'\) contains \(p_i\) or \(\neg p_i\), so
  \(\mathcal C'\) contains at most \(n-1\) distinct symbols.

  \emph{Claim: \(\mathcal C'\) is unsatisfiable.}
  Suppose instead that \(v(\mathcal C')=1\). We show that changing only
  the value of \(p_i\), if necessary, gives a satisfying assignment for
  \(\mathcal C\). There are two cases.
  \begin{enumerate}
    \item Suppose \(v(C\setminus\{\neg p_i\})=1\) for every
      \(C\in\mathcal C_{01}\). Define \(w(p_i)=1\) and
      \(w(p_j)=v(p_j)\) for \(j\ne i\). Each clause of
      \(\mathcal C_{01}\) remains true because its part not involving
      \(p_i\) is true. Each clause of \(\mathcal C_{10}\) is true because
      it contains \(p_i\).
    \item Otherwise, choose \(C_1\in\mathcal C_{01}\) such that
      \(v(C_1\setminus\{\neg p_i\})=0\). For every
      \(C_2\in\mathcal C_{10}\), the resolvent
      \[
        (C_1\setminus\{\neg p_i\})\cup(C_2\setminus\{p_i\})
      \]
      belongs to \(\mathcal C'\) and is true under \(v\). Its first part
      is false, so \(v(C_2\setminus\{p_i\})=1\).
      Define \(w(p_i)=0\) and \(w(p_j)=v(p_j)\) for \(j\ne i\).
      Every clause of \(\mathcal C_{01}\) is true because it contains
      \(\neg p_i\); every clause of \(\mathcal C_{10}\) is true because
      its remaining part is true.
  \end{enumerate}
  In both cases, the clauses of \(\mathcal C_{00}\) remain true because
  their truth values do not depend on \(p_i\), and all clauses of
  \(\mathcal C_{11}\) are true under every assignment. Thus
  \(w(\mathcal C)=1\), contradicting the unsatisfiability of \(\mathcal C\).
  This proves the claim.

  By the induction hypothesis, \(\mathcal C'\vdash_R\emptyset\).
  To obtain a proof from \(\mathcal C\), first list the clauses of
  \(\mathcal C\) as assumptions and derive all clauses in
  \(\mathcal C_{01}\times_R\mathcal C_{10}\) by resolution. Every assumption
  used in the refutation from \(\mathcal C'\) is now available: it is either
  in \(\mathcal C_{00}\subseteq\mathcal C\) or among these resolvents.
  Reproduce the resolution steps of that refutation, referring to the
  already available lines whenever it uses an assumption. This yields a
  finite resolution proof of \(\emptyset\) from \(\mathcal C\).
\end{proof}

\begin{remark}[Scope of the proof]
  The course book states Lemma 6.7 for arbitrary clause sets, but the
  induction above proves the finite case treated in the lecture. The
  general case follows using propositional compactness: every unsatisfiable
  clause set has a finite unsatisfiable subset, whose refutation is also a
  refutation from the original set. Compactness is an additional result;
  it is not proved by this induction.
\end{remark}

\begin{corollary}[Refutation criterion]
  For a finite clause set \(\mathcal C\),
  \[
    \mathcal C\text{ is unsatisfiable}
    \quad\Longleftrightarrow\quad
    \mathcal C\vdash_R\emptyset.
  \]
\end{corollary}

\begin{proof}
  The forward direction is refutation completeness. Conversely, if
  \(\mathcal C\vdash_R\emptyset\), soundness says that any assignment
  satisfying \(\mathcal C\) must satisfy \(\emptyset\). Since the empty
  clause is always false, no such assignment exists.
\end{proof}

\end{document}
